\documentclass[11pt,a4paper]{amsart}
\usepackage[T1]{fontenc}
\usepackage{lmodern}
\usepackage[margin=25mm]{geometry}
\usepackage{amsmath,amssymb,amsthm}
\usepackage[expansion=false]{microtype}
\usepackage{array,booktabs,tabularx}
\usepackage{needspace}
\setlength{\captionindent}{0pt}
\usepackage{xurl}
\usepackage[hidelinks]{hyperref}
\urlstyle{same}
\hypersetup{
  pdftitle={Erd\H{o}s Problem 928: quantitative comparisons and original formulations},
  pdfauthor={JI HO BAE},
  pdfsubject={Elementary comparison estimates and the deduction of Erd\H{o}s Problem 928 from OpenAI's joint Dickman law}
}
\newtheorem{proposition}{Proposition}
\newtheorem{corollary}[proposition]{Corollary}
\newcommand{\N}{\mathbb N}
\newcommand{\A}{\mathcal A}
\newcommand{\B}{\mathcal B}
\newcommand{\C}{\mathcal C}
\newcolumntype{L}[1]{>{\raggedright\arraybackslash}p{#1}}
\newcolumntype{Y}{>{\raggedright\arraybackslash}X}
\title[Erd\H{o}s Problem 928 and the joint Dickman law]{Erd\H{o}s Problem 928:\\
quantitative comparisons and original formulations}
\author{JI HO BAE}
\date{October 9, 2026}

\begin{document}
\begin{abstract}
We prove explicitly that OpenAI's joint Dickman law implies the density
formula $\rho(1/\alpha)\rho(1/\beta)$ in the precise formulation of
Erd\H{o}s Problem~928. Here $\rho$ is the Dickman--de Bruijn function.
For fixed $0<\alpha,\beta<1$ and real $X\ge2$, an elementary argument
bounds the symmetric difference between the counting sets in the problem
and in the cited theorem, restricted to $2\le n\le X$, by
$\pi(X^\alpha)+\pi((X+1)^\beta)$, where $\pi$ counts primes.
This bound handles both the strict inequalities and the change in the
second threshold from $n^\beta$ to $(n+1)^\beta$. We also treat the
strict formulation with both thresholds based at $n$, and compare the
formulations in Erd\H{o}s's two 1976 papers and the 1978
Erd\H{o}s--Pomerance paper. The analytic joint law is due to OpenAI;
our contribution consists of the explicit finite comparison estimates
and the precise correspondence among these formulations.
\end{abstract}
\maketitle

\section{Introduction}

For an integer $n\ge2$, let $P(n)$ be the largest prime factor of $n$, and set
$P(1)=1$. Write $\rho$ for the Dickman--de Bruijn function: the continuous
function on $[0,\infty)$ determined by
\[
 \rho(u)=1\quad(0\le u\le1),\qquad
 u\rho'(u)=-\rho(u-1)\quad(u>1).
\]
The classical smooth-number limit is
\begin{equation}\label{eq:marginal}
 \lim_{X\to\infty}\frac1X
 \#\{1\le n\le X:P(n)\le X^c\}=\rho(1/c)
 \qquad(0<c<1);
\end{equation}
see de Bruijn~\cite[p.~51, (1.1)--(1.2)]{deBruijn1951}.
We also use the standard properties $\rho(u)>0$ for finite $u\ge0$
and $\rho(u)\to0$ as $u\to\infty$; see
\cite[p.~52, (1.8)]{deBruijn1951} together with the defining equation.
We use $\N=\{1,2,\ldots\}$. All counting variables are integers, and
the endpoints $X$ are real. Natural density means the limit of the
ordinary, unweighted count up to $X$, divided by $X$.

Problem~928 in Bloom's collection~\cite{Bloom928} asks whether, for fixed
$\alpha,\beta\in(0,1)$, the integers satisfying
\begin{equation}\label{eq:problem}
 P(n)<n^\alpha,\qquad P(n+1)<(n+1)^\beta
\end{equation}
have a natural density. The accompanying discussion asks whether this
density is $\rho(1/\alpha)\rho(1/\beta)$. The page cites
Erd\H{o}s~\cite{Erdos1976Manitoba,Erdos1976Debrecen} and
Erd\H{o}s--Pomerance~\cite{ErdosPomerance1978}; their formulations will
be compared in Section~\ref{sec:sources}.

Our analytic input is OpenAI's joint Dickman
law~\cite[Theorem~1.1, p.~2]{OpenAI2026}: for every fixed
$\alpha,\beta\in(0,1)$,
\begin{equation}\label{eq:JD}
 \lim_{X\to\infty}\frac1X
 \#\{2\le n\le X:P(n)\le n^\alpha,\ P(n+1)\le n^\beta\}
 =\rho(1/\alpha)\rho(1/\beta).
\end{equation}
OpenAI identifies its result as resolving the
Erd\H{o}s--Pomerance joint Dickman conjecture. Section~11.3 of that
paper also proves the versions with thresholds $X^\alpha,X^\beta$ and
the corresponding upper-tail conditions. We take these analytic results
as given. Proposition~\ref{prop:comparison} supplies explicit elementary
estimates for the changes of convention needed in~\eqref{eq:problem}
and in Erd\H{o}s's strict formulation. Section~\ref{sec:sources} identifies
the precise original statements to which the results apply.

\section{Finite comparison estimates}\label{sec:comparison}

Fix $\alpha,\beta\in(0,1)$. For $X\ge2$, define
\begin{align*}
 \A_{\alpha,\beta}(X)
   &=\{n\in\N:2\le n\le X,\ P(n)<n^\alpha,
                                \ P(n+1)<(n+1)^\beta\},\\
 \B_{\alpha,\beta}(X)
   &=\{n\in\N:2\le n\le X,\ P(n)\le n^\alpha,
                                \ P(n+1)\le n^\beta\},\\
 \C_{\alpha,\beta}(X)
   &=\{n\in\N:2\le n\le X,\ P(n)<n^\alpha,
                                \ P(n+1)<n^\beta\}.
\end{align*}
The sets $\A$ and $\C$ correspond respectively to Problem~928 and
condition~(1) in~\cite[p.~25]{Erdos1976Manitoba}; $\B$ is the set in
\eqref{eq:JD}. Let $\pi(y)=\#\{p\le y:p\text{ is prime}\}$, and let
$S\mathbin\triangle T$ denote the symmetric difference of two sets.

\begin{proposition}[Comparison of conventions]\label{prop:comparison}
For every $\alpha,\beta\in(0,1)$ and every real $X\ge2$,
\begin{align}
 \#\bigl(\A_{\alpha,\beta}(X)\mathbin\triangle
          \B_{\alpha,\beta}(X)\bigr)
 &\le \pi(X^\alpha)+\pi((X+1)^\beta),\label{eq:AB}\\
 0\le\#\B_{\alpha,\beta}(X)-\#\C_{\alpha,\beta}(X)
 &\le \pi(X^\alpha)+\pi(X^\beta).\label{eq:BC}
\end{align}
\end{proposition}

\begin{proof}
Suppress the subscripts and the argument $X$. If $n\in\B\setminus\A$,
then $P(n+1)\le n^\beta<(n+1)^\beta$, so the second condition for $\A$
holds. The first condition must therefore fail, giving $P(n)=n^\alpha$.
The prime $p=P(n)\le X^\alpha$ determines $n=p^{1/\alpha}$ uniquely.
Consequently $\#(\B\setminus\A)\le\pi(X^\alpha)$.

If $n\in\A\setminus\B$, the first condition for $\B$ holds, so
\[
 n^\beta<P(n+1)<(n+1)^\beta.
\]
The intervals $(m^\beta,(m+1)^\beta)$, for integers $m\ge2$, are
pairwise disjoint because $\beta>0$. Thus $n\mapsto P(n+1)$ is injective
on $\A\setminus\B$, and its values are primes less than $(X+1)^\beta$.
This proves $\#(\A\setminus\B)\le\pi((X+1)^\beta)$ and hence~\eqref{eq:AB}.

Finally, $\C\subseteq\B$. Any $n\in\B\setminus\C$ satisfies
$P(n)=n^\alpha$ or $P(n+1)=n^\beta$. The first equality contributes at
most $\pi(X^\alpha)$ integers, by the argument above. For the second,
$n=P(n+1)^{1/\beta}$ is determined by a prime at most $X^\beta$,
giving at most $\pi(X^\beta)$ integers. Summing these two bounds
proves~\eqref{eq:BC}.
\end{proof}

\begin{corollary}[Problem~928 and Erd\H{o}s's strict formulation]\label{cor:density}
For every fixed $\alpha,\beta\in(0,1)$,
\begin{equation}\label{eq:conclusion}
 \lim_{X\to\infty}\frac{\#\A_{\alpha,\beta}(X)}X
 =\lim_{X\to\infty}\frac{\#\C_{\alpha,\beta}(X)}X
 =\rho(1/\alpha)\rho(1/\beta).
\end{equation}
In particular, the density in Problem~928 and the density
$d(\alpha,\beta)$ asked for in~\cite[p.~25]{Erdos1976Manitoba} both
exist and equal the predicted product.
\end{corollary}

\begin{proof}
For finite sets, $|\#S-\#T|\le\#(S\mathbin\triangle T)$. Since
$\pi(y)\le\lfloor y\rfloor$, both upper bounds in
Proposition~\ref{prop:comparison} are $O(X^\alpha+X^\beta)=o(X)$.
After division by $X$, equation~\eqref{eq:JD} therefore gives both
limits in~\eqref{eq:conclusion}.
\end{proof}

The comparison estimates are independent of the joint Dickman law.
They quantify the change of convention for fixed $\alpha,\beta$;
the asymptotic evaluation of the count is supplied by~\eqref{eq:JD}.

\section{Comparison with the original sources}\label{sec:sources}

We call a threshold \emph{moving} when it depends on the integer being
counted, and \emph{fixed} when it is held constant over the counting
range. Thus $n^\alpha$ and $(n+1)^\beta$ are moving thresholds, while
$X^\alpha$ and $X^\beta$ are fixed thresholds. Put
\[
 a(0)=1,\qquad a(t)=1-\rho(1/t)\quad(0<t\le1).
\]
Table~\ref{tab:forms} summarizes the two-integer formulations. The
subsections below justify the changes of convention and counting range.

\begin{table}[htbp]
\small
\caption{The two-integer formulations and their deductions.
Exponents are fixed: \mbox{$0<\alpha,\beta<1$},
\mbox{$\gamma>1$}, and \mbox{$0\le t,s\le1$}.
Unless indicated otherwise, \mbox{$2\le n\le X$}.}
\label{tab:forms}
\begin{tabularx}{\textwidth}{@{}L{32mm}Y L{46mm}@{}}
\toprule
Source & Conditions & Deduction\\
\midrule
\cite{Bloom928}; independence statement in
\cite[p.~25]{Erdos1976Manitoba}
 & $P(n)<n^\alpha$, $P(n+1)<(n+1)^\beta$
 & Corollary~\ref{cor:density}, set $\A$: density
   $\rho(1/\alpha)\rho(1/\beta)$\\[3pt]
\cite[p.~25, (1)]{Erdos1976Manitoba}
 & $P(n)<n^\alpha$, $P(n+1)<n^\beta$
 & Corollary~\ref{cor:density}, set $\C$: density
   $\rho(1/\alpha)\rho(1/\beta)$\\[3pt]
\cite[pp.~272--273]{Erdos1976Debrecen}, $r=2$
 & $P(n+1)\le k$, $P(n+2)\le k$; $1\le n<k^\gamma$
 & OpenAI~\cite[(11.7), p.~80]{OpenAI2026}:
   count $\sim\rho(\gamma)^2 k^\gamma$\\[3pt]
\cite[p.~311]{ErdosPomerance1978}
 & $P(n)\ge X^t$, $P(n+1)\ge X^s$
 & OpenAI~\cite[(11.8) and following paragraph, p.~80]{OpenAI2026}:
   count$/X\to a(t)a(s)$\\
\bottomrule
\end{tabularx}
\end{table}

\subsection{Erd\H{o}s's Manitoba paper}
On page~25 of~\cite{Erdos1976Manitoba}, Erd\H{o}s expects independence
for $P(n)<n^\alpha$ and $P(n+1)<(n+1)^\beta$. In the following sentence,
he asks whether the integers satisfying his condition~(1), namely
$P(n)<n^\alpha$ and $P(n+1)<n^\beta$, have a density $d(\alpha,\beta)$.
Thus the webpage uses the inequalities in his independence statement,
whereas condition~(1) uses $n^\beta$ for the second threshold.
Corollary~\ref{cor:density} treats both forms explicitly.

Erd\H{o}s also asks whether these strict inequalities have infinitely
many solutions for each fixed pair $\alpha,\beta>0$. The positive value in
\eqref{eq:conclusion} implies this, also for parameters outside $(0,1)$
by choosing smaller positive parameters inside that interval.
Infinitude itself was already available: as recorded in~\cite{Bloom928},
it follows from Schinzel~\cite[Theorem~14, p.~231]{Schinzel1967}, applied
to $(2n+1)^2-1=4n(n+1)$. It gives a constant $K>0$ and infinitely many
integers $n\ge3$ such that
\[
 P(n(n+1))\le n^{K/\log\log n}.
\]
Since $K/\log\log n\to0$, these integers satisfy Erd\H{o}s's strict
inequalities for every prescribed positive pair $\alpha,\beta$, once
$n$ is sufficiently large.

\subsection{Erd\H{o}s's Debrecen paper}
Write $\gamma>1$ for the exponent denoted by $\alpha$
in~\cite[pp.~272--273]{Erdos1976Debrecen}. Erd\H{o}s recalls the
one-variable asymptotic $(c_\gamma+o(1))k^\gamma$ for the number of
$m\le k^\gamma$ with $P(m)\le k$, and conjectures, for each fixed integer
$r\ge1$, that as $k$ tends to infinity through positive integers,
\begin{equation}\label{eq:Debrecen}
 \#\{1\le n<k^\gamma:P((n+1)\cdots(n+r))\le k\}
  =(c_\gamma^r+o(1))k^\gamma.
\end{equation}
Equation~\eqref{eq:marginal} identifies $c_\gamma=\rho(\gamma)$.

OpenAI~\cite[(11.7), p.~80]{OpenAI2026} proves, for fixed $c,d\in(0,1)$,
\begin{equation}\label{eq:fixed}
 \lim_{X\to\infty}\frac1X
 \#\{2\le m\le X:P(m)\le X^c,\ P(m+1)\le X^d\}
 =\rho(1/c)\rho(1/d).
\end{equation}
For $r=2$ in~\eqref{eq:Debrecen}, put $X=k^\gamma$, $m=n+1$, and
$c=d=1/\gamma$. Since $P(m(m+1))=\max\{P(m),P(m+1)\}$, the
prime-factor conditions coincide with those in~\eqref{eq:fixed}.
The ranges $1\le n<X$ and $2\le m\le X$ differ,
after the shift, by at most one endpoint term. Hence~\eqref{eq:fixed}
gives~\eqref{eq:Debrecen} for $r=2$.

The webpage~\cite{Bloom928} states the higher-order conjecture using
$P(n+i)\le n^\alpha$ for $0\le i<r$. Its case $r=2$ is directly
\eqref{eq:JD} with $\beta=\alpha$. The results used here concern two
consecutive integers and do not establish the $r\ge3$ assertion.

\Needspace{10\baselineskip}
\subsection{Erd\H{o}s and Pomerance}
On page~311 of~\cite{ErdosPomerance1978}, the authors use fixed
thresholds and upper tails. For $0\le t,s\le1$, they ask whether
\begin{equation}\label{eq:upper}
 b(t,s)=\lim_{X\to\infty}\frac1X
 \#\{2\le n\le X:P(n)\ge X^t,\ P(n+1)\ge X^s\}
\end{equation}
exists, is continuous on $[0,1]^2$, and equals $a(t)a(s)$, where their
$a(t)$ is the corresponding one-variable upper-tail limit.
Equation~\eqref{eq:marginal}, together with the treatment of equality and
endpoints below, identifies their $a(t)$ with the function defined above.

For $0<t,s<1$, inclusion--exclusion applied to~\eqref{eq:fixed} and the
two marginal laws gives the limit $a(t)a(s)$ with strict upper tails
$P(n)>X^t$, $P(n+1)>X^s$. For the shifted marginal, replacing $n$ by
$n+1$ changes the counting range by only $O(1)$ terms.

To pass to weak upper tails, observe that for fixed $0<t\le1$ the
equality $P(m)=X^t$ is
impossible unless $X^t$ is a prime $p$. In that case $p\mid m$, so
there are at most $\lfloor(X+1)/p\rfloor=O(X^{1-t}+1)=o(X)$ such
$m\le X+1$.

At exponent $0$ the upper-tail condition is automatic for $m\ge2$.
At exponent $1$, within the range $2\le m\le X+1$, the condition
$P(m)\ge X$ forces $m\in[X,X+1]$, so at most two integers occur.
This proves~\eqref{eq:upper} at all
endpoints as well.
Since $\rho(1)=1$ and $\rho(u)\to0$ as $u\to\infty$, the function $a$
is continuous on $[0,1]$, and so is $b(t,s)=a(t)a(s)$.

This deduction, including the treatment of equality and endpoints, is
already given in~\cite[(11.8) and following paragraph, p.~80]{OpenAI2026}.
The role of the present comparison is to match it to the exact statement
in~\cite{ErdosPomerance1978}. All the two-integer formulations in
Table~\ref{tab:forms} therefore follow from the cited joint Dickman law
and its fixed-threshold version.

\begin{thebibliography}{9}
\bibitem{Bloom928}
T.~F. Bloom, \emph{Erd\H{o}s Problem \#928},
accessed October~9, 2026.
\newline\url{https://www.erdosproblems.com/928}.

\bibitem{deBruijn1951}
N.~G. de Bruijn, \emph{On the number of positive integers $\le x$ and
free of prime factors $>y$}, Nederl. Akad. Wetensch. Proc. Ser. A
\textbf{54} (1951), 50--60.
\newline\url{https://pure.tue.nl/ws/files/1793853/597499.pdf}.

\bibitem{Erdos1976Manitoba}
P.~Erd\H{o}s, \emph{Problems and results on number theoretic properties of
consecutive integers and related questions}, Proceedings of the Fifth
Manitoba Conference on Numerical Mathematics (Winnipeg, 1975), Congressus
Numerantium XVI, Utilitas Math., Winnipeg, 1976, pp.~25--44.
\newline\url{https://users.renyi.hu/~p_erdos/1976-39.pdf}.

\bibitem{Erdos1976Debrecen}
P.~Erd\H{o}s, \emph{Problems and results on consecutive integers}, Publ.
Math. Debrecen \textbf{23} (1976), 271--282.
\newline\url{https://doi.org/10.5486/PMD.1976.23.3-4.15}.

\bibitem{ErdosPomerance1978}
P.~Erd\H{o}s and C.~Pomerance, \emph{On the largest prime factors of $n$
and $n+1$}, Aequationes Math. \textbf{17} (1978), 311--321.
\newline\url{https://users.renyi.hu/~p_erdos/1978-29.pdf}.

\bibitem{OpenAI2026}
OpenAI, \emph{The joint Dickman law for consecutive integers}, OpenAI
Math Release preprint, September~24, 2026, commit
\texttt{fd4aeeb2ee4f}.
\newline\url{https://github.com/openai/math/blob/fd4aeeb2ee4fc729c18d98444fed42fd0529eeeb/preprints/The-joint-Dickman-law-for-consecutive-integers-September-24-2026/paper.pdf}.

\bibitem{Schinzel1967}
A.~Schinzel, \emph{On two theorems of Gelfond and some of their
applications}, Acta Arith. \textbf{13} (1967), 177--236.
\newline\url{https://doi.org/10.4064/aa-13-2-177-236}.
\end{thebibliography}

\end{document}
